\section{Research-to-Experiment Traceability}
The study's single two-phase conceptual framework is presented earlier in Chapter~3. For the completed experimental phase, researchers selected scenario categories, generated a controlled corpus, partitioned related scenarios before fitting text features, tuned the baseline and SVM on the training partition, and evaluated the fixed models on separately held-out partitions. These steps are explained here without repeating the conceptual framework diagram. The single combined evaluation-protocol figure is Figure~\ref{fig:evaluation}.\par

\section{Data Source and Scenario Acquisition Protocol}
The project does not claim to have collected 24,000 authentic resident reports. \emph{Data acquisition} refers here to acquiring a documented \emph{evidence foundation} for generating fictional, text-only scenarios. The category register has 60 report-topic entries, 28 clarification/review-reason entries, and 30 spam/non-complaint entries. The source-to-category matrix links each entry to a suitable law, administrative reporting practice, linguistic study, or operational labeling rule. For example, the Local Government Code \cite{ra7160} provides the broad basis for public-service reporting; RA 9003 \cite{ra9003} provides subject-specific environmental context; Davao City Reports \cite{dcr} illustrates expected descriptions of location and incident details; and the Filipino complaint-expression studies \cite{deleon,quimbo} support alternatives to rigid report templates. An authority that establishes the existence of a reporting topic does \emph{not} establish the relative frequency or authenticity of each synthetic narrative.

\begin{table}[H]\centering\caption{Documented roles of evidence used in controlled generation}\small\begin{tabularx}{\textwidth}{p{.22\textwidth}Xp{.23\textwidth}}\toprule
Evidence class & Generation decision supported & Limit of inference\\\midrule
Philippine legislation & Legal or administrative relevance of a public concern & Does not define the three model labels\\
Official complaint workflows & Common incident fields and missing-context issues & Cannot prove a given report is false\\
Complaint-pragmatics research & Directness, criticism, questions, requests and affect & Does not provide nationwide style proportions\\
Code-switching research & Plausible multilingual phrase ordering and lexical borrowing & Not direct evidence of barangay frequencies\\
Scam and spam studies & Promotional requests, impersonation and solicitations & Does not establish all unrelated-message types\\
Synthetic data research & Diversity control, documentation and independent evaluation & Does not guarantee field generalization\\\bottomrule
\end{tabularx}\end{table}

\section{Labeling Codebook and Error-Sensitive Boundaries}
The primary label is operational and text-based: \textbf{Valid} indicates an understandable report of a civic issue or a request for action, regardless of whether the barangay itself has jurisdiction. \textbf{Suspicious} indicates that a plausible report cannot be interpreted sufficiently because essential meaning is missing or internally contradictory. \textbf{Spam} indicates an irrelevant, deceptive, promotional, or non-complaint submission. These labels do not establish truthfulness or determine legal jurisdiction. A brief Cebuano report with a clear problem is not Suspicious because of brevity alone; a complaint about someone attempting a scam is not Spam merely because it contains scam-related vocabulary. The codebook and category-level evidence are supplied with the dataset.

If a condition depends on information the classifier never receives (for example, an absent file attachment or a separate form field), the system should assign the check to application validation or a human reviewer rather than learning an inference unsupported by \texttt{complaint\_text}. Urgent public-safety and protection-related cases must not be blocked solely by a classifier label. This is a design safeguard, not an evaluated software behavior in the present study.

\section{Two-Level Validation Design}
For the present study, the experiment uses \textbf{two distinct grouped resampling procedures}. First, a three-fold \texttt{GridSearchCV} searches 18 SVM configurations on the 19,213 training samples. Second, an additional five-fold \texttt{StratifiedGroupKFold} is run on the same training partition for the fixed chosen SVM and Naive Bayes baseline to report fold-specific stability. The groups are \texttt{scenario\_family\_id} identifiers, not the three class labels, and each held-out fold has zero overlapping group IDs with its corresponding training fold. The vectorizer is fitted anew inside every fold. The held-out validation set comprises 2,399 examples and is used for validation metrics, class confusion, language subgroup comparisons, and descriptive feature-ablation diagnostics. The 2,388-record final test set was excluded from the supplementary five-fold suite; it had already been evaluated in the original full-results experiment and is not an untouched future test set.

\section{Evaluation Score Reporting and Fold Dispersion}
The primary scoring metric is Macro-F1; Chapter~3's Milestone 2 evaluation-metrics section defines accuracy, class precision, recall, F1, macro- and weighted-F1, and the confusion matrix. For the separate five-fold stability analysis, the arithmetic fold mean is $\bar{F}=\frac15\sum_{i=1}^5F_i$ and the sample standard deviation is $s=\sqrt{\sum_i(F_i-\bar F)^2/4}$. These describe dispersion across generated-data folds; they are not independent population confidence intervals.\par

\section{Reproducibility and Validation Integrity}
The accompanying grouped-validation script includes assertions for the 24,000-record count, training/validation group separation, and zero within-fold group overlap. It exports the fold composition, five-fold metrics, validation predictions, per-class and language-subgroup metrics, confusion matrices, and diagnostic ablations. The recorded seed is \texttt{20261026}, and the environment uses scikit-learn version 1.8.0. The full experiment still depends on the stated generator and label-rubric assumptions. The additional recorded diagnostics must be interpreted as a retrospective validation supplement and not as a preregistered independent test.
